\section{Metodología de solución}
\subsection{Selección de zonas candidatas}
Se ordenan las celdas por probabilidad y se exige separación mínima entre candidatos para evitar múltiples destinos equivalentes dentro de la misma mancha espacial.
\subsection{Ruteo con A*}
La grilla válida se transforma en un grafo de navegación. A* busca la ruta de costo mínimo utilizando distancia Haversine. Si existe una máscara legal, los nodos prohibidos se eliminan antes del ruteo.
\subsection{Greedy}
Asigna secuencialmente cada barco a la alternativa con mayor puntuación inmediata, incluyendo penalización por distancia y congestión.
\subsection{MILP}
Optimiza de forma simultánea la asignación completa de la flota y funciona como benchmark centralizado para el escenario estático.
\subsection{MARL}
Se implementa una política MARL cooperativa con parámetros compartidos y actualización mediante gradiente de política. El entrenamiento utiliza directamente las zonas y probabilidades obtenidas de \texttt{MapProbabilidad\_adulto.csv}; no se genera ni se perturba una superficie sintética de probabilidad. La estocasticidad corresponde únicamente al muestreo de acciones durante el aprendizaje.
\subsection{Flujo técnico}
\begin{figure}[H]
\centering
\begin{tikzpicture}[node distance=6mm and 5mm,box/.style={rectangle,rounded corners,draw=navy,thick,fill=softblue,align=center,minimum height=9mm,text width=3cm},arrow/.style={-{Latex[length=2mm]},thick,draw=navy}]
\node[box] (a) {Mapa de\\probabilidad};
\node[box,right=of a] (b) {Máscara legal\\SERNANP};
\node[box,right=of b] (c) {Grafo\\navegable};
\node[box,below=of c] (d) {Zonas\\candidatas};
\node[box,left=of d] (e) {Rutas A*};
\node[box,left=of e] (f) {Greedy / MILP /\\MARL};
\node[box,below=of f] (g) {Mapas, costos y\\simulación};
\draw[arrow] (a)--(b);\draw[arrow] (b)--(c);\draw[arrow] (c)--(d);\draw[arrow] (d)--(e);\draw[arrow] (e)--(f);\draw[arrow] (f)--(g);
\end{tikzpicture}
\caption{Arquitectura general del sistema.}
\end{figure}